% !TEX root = ../main.tex
%% Devoir surveillé n° 1 — chapitres 1 à 4 (valeurs vérifiées : verifications/devoirs.py)

\devoirDS{1}{Chapitres 1 à 4}{{Le bleu d'un sirop de menthe : dosage par étalonnage}/8,{Argenter un fil de cuivre : oxydoréduction et avancement}/7,{Une boisson sucrée : quantité de matière et concentration}/5}

\begin{exoBac}{8}{Le bleu d'un sirop de menthe}
Certains sirops de menthe doivent leur couleur bleue à un colorant alimentaire, le bleu patenté V (E131). On souhaite
vérifier que la consommation de ce sirop ne présente pas de risque pour un enfant.
\begin{donnees}
Dose journalière admissible (DJA) du bleu patenté V : \qty{5}{mg} par kilogramme de masse corporelle et par jour. Verrerie
disponible : pipettes jaugées de 5,0 ; 10,0 et \qty{20,0}{mL} ; fioles jaugées de 50,0 ; 100,0 et \qty{250,0}{mL} ; béchers ; pipette
simple ; eau distillée.
\end{donnees}
\begin{docu}{1}{Spectre d'absorption du bleu patenté V}
\schema{ds1-spectre}
\end{docu}
\begin{docu}{2}{Étalonnage et mesure}
On mesure l'absorbance de solutions étalons de bleu patenté V à \qty{638}{nm}. Le sirop, dilué 20 fois, a une absorbance $A=0{,}52$.\par\smallskip
{\centering\small\begin{tabular}{@{}lccccc@{}}\toprule
$c$ (\unit{mg.L^{-1}}) & 1,0 & 2,0 & 3,0 & 4,0 & 5,0 \\
$A$ & 0,16 & 0,33 & 0,48 & 0,64 & 0,81 \\\bottomrule
\end{tabular}\par}
\schema{ds1-etalonnage}
\end{docu}
\begin{questions}
\item Expliquer pourquoi une solution de bleu patenté V paraît bleue.
\item Justifier le choix de la longueur d'onde de \qty{638}{nm} pour les mesures.
\item Montrer que les solutions étalons vérifient la loi de Beer-Lambert et déterminer le coefficient $k$ de la relation $A=k\,c$.
\item Choisir la verrerie nécessaire pour diluer 20 fois le sirop et décrire le protocole.
\item Déterminer la concentration en masse de bleu patenté V dans la solution diluée, puis dans le sirop.
\item \initiative La boisson est préparée avec un volume de sirop pour sept volumes d'eau. Combien de verres de \qty{25}{cL} un
enfant de \qty{30}{kg} devrait-il boire dans la journée pour atteindre la DJA ? Conclure.
\end{questions}
\end{exoBac}

\begin{exoBac}{7}{Argenter un fil de cuivre}
On plonge un fil de cuivre de masse $m=\qty{0,635}{g}$ dans $V=\qty{50,0}{mL}$ d'une solution de nitrate d'argent, $\ce{Ag+}+\ce{NO3-}$, de
concentration $c=\qty{0,100}{mol.L^{-1}}$. Après quelques heures, le fil s'est recouvert d'un dépôt d'argent métallique et la solution
est devenue bleue.
\begin{donnees}
Masses molaires : $M(\ce{Cu})=\qty{63,5}{g.mol^{-1}}$ ; $M(\ce{Ag})=\qty{107,9}{g.mol^{-1}}$. Couples oxydant-réducteur : $\ce{Ag+}/\ce{Ag}$ et
$\ce{Cu^{2+}}/\ce{Cu}$. En solution, les ions $\ce{Cu^{2+}}$ sont bleus ; les ions $\ce{Ag+}$ et $\ce{NO3-}$ sont incolores.
\end{donnees}
\begin{questions}
\item Écrire les demi-équations électroniques des deux couples. Identifier l'oxydant et le réducteur qui réagissent.
\item En déduire l'équation de la réaction.
\item Calculer les quantités de matière initiales des réactifs.
\item Construire le tableau d'avancement. Déterminer l'avancement maximal et le réactif limitant.
\item Déterminer la masse de cuivre restante, la masse d'argent déposée et la concentration finale en ions $\ce{Cu^{2+}}$.
\item On sort le fil, recouvert de son dépôt, et on le sèche. Quelle masse doit indiquer la balance ?
\end{questions}
\end{exoBac}

\begin{exoBac}{5}{Une boisson sucrée}
Une canette de \qty{33}{cL} de soda contient \qty{35}{g} de sucre, le saccharose, de formule brute $\ce{C12H22O11}$.
\begin{donnees}
Masses molaires atomiques : $M(\ce{C})=\qty{12,0}{g.mol^{-1}}$ ; $M(\ce{H})=\qty{1,0}{g.mol^{-1}}$ ; $M(\ce{O})=\qty{16,0}{g.mol^{-1}}$. Constante
d'Avogadro : $N_\mathrm{A}=\qty{6,02e23}{mol^{-1}}$.
\end{donnees}
\begin{questions}
\item Calculer la masse molaire du saccharose.
\item Calculer la concentration en masse, puis la concentration en quantité de matière de saccharose dans le soda.
\item Combien de molécules de saccharose la canette contient-elle ?
\item Décrire la préparation de \qty{100,0}{mL} d'une solution de saccharose de même concentration que le soda, à partir de
saccharose solide.
\item Quel volume de cette solution faut-il prélever pour préparer \qty{50,0}{mL} d'une solution dix fois moins concentrée ?
\end{questions}
\end{exoBac}
